\documentclass[11pt]{article}
\usepackage[margin=1in]{geometry}
\usepackage{amsmath,amssymb}
\usepackage[hidelinks]{hyperref}

\begin{document}

\section*{Human Review: nn-Bounded Operators and nn-Convergence}

\noindent Original automatic proof: \texttt{solution\_packet.pdf}

\noindent Checked date: 15 June 2026

\noindent Status: Rejected as literature already answered.

\noindent Verdict: This packet should not be counted as a novel counterexample. The question was answered in later literature, specifically in Theorem 2.2 of Zabeti's paper \emph{Topological Algebras of Bounded Operators on Topological Vector Spaces}.

\section*{Human Comments}

The automatic packet addresses Troitsky's question asking whether the class of nn-bounded operators is closed under nn-convergence. The proposed construction uses finite truncations of the left shift on \(\mathbb R^{\mathbb N}\), with the product topology, converging in nn-convergence to the full left shift, which is not nn-bounded.

The mathematical construction is valid, but it should be excluded from the survey as a new discovery, because the same phenomenon is already treated in the literature:
\[
\text{Zabeti, \emph{Topological Algebras of Bounded Operators on Topological Vector Spaces},}
\]
\[
\text{\emph{Journal of Advanced Research in Pure Mathematics} 3(1) (2011), 22--26.}
\]
The DOI metadata for this paper is
\[
\texttt{10.5373/jarpm.438.052210}.
\]

More specifically, Theorem 2.2 of that paper states that, for a locally pseudo-convex topological vector space \(X\), the operations in \(B_{nn}(X)\) are continuous for the topology of nn-convergence, and that \(B_{nn}(X)\) fails to be closed in this topology. The proof gives the same kind of example: take \(X=\mathbb C^{\mathbb N}\) with the topology of coordinatewise convergence, let \(S_n\) be the left shift followed by projection onto the first \(n\) coordinates, and observe that \(S_n\) nn-converges to the full left shift \(S\), while \(S\) is not nn-bounded, as in Troitsky's Example 2.5.

Thus this packet is best classified as \emph{already answered by later work}, rather than as a new counterexample produced by the pipeline. It appears to be a rediscovery of the example in Zabeti's Theorem 2.2. I would therefore reject it for inclusion among novel or potentially novel mathematical outputs.

\section*{Review Outcome}

Rejected. The reason is not that the construction is incorrect, but that the question was already answered in the later literature cited above. This should be moved conceptually to the literature-already-answered category if the repository classification is updated later.

\end{document}
